\scanpage{17}
\uline{Si} $n+\rho$ \uline{est impair}, alors $b_n=\dim\underbrace{H^n(\bar X,\bar E)}_{\text{poids }n+\rho}$ est pair, donc $\deleted{\unsure{\varepsilon}}(-1)^{b_n}=1$. De plus, un raisonnement du cours de \unsure{Serre} utilisant une forme alternée non dégénérée \unsure{invariante} montre que dans ce cas $\operatorname{signe}\det f_{H^n(\bar X,\bar E)}$ est égal à $+1$. Donc on trouve alors $\varepsilon'(E)=1$.

Si $n+\rho$ est pair, on ne peut déterminer le signe de façon absolue. On peut \deleted{\unsure{seulement affirmer que}} \margin{(\unsure{analogue de la parité|anal. de la parité})}
si $\unsure{\varepsilon_1|e_1}$ ($\unsure{\varepsilon_{-1}|e_{-1}}$) est \unsure{la} \deleted{\unsure{dimension}} \inserted{multiplicité} de \deleted{\unsure{l'espace}} la valeur propre $1$ ($-1$) de $f$ sur $H^n(\bar X,\bar F)$, on a
\[
(-1)^{b_n}=(-1)^{\unsure{\varepsilon_1+\varepsilon_{-1}|e_1+e_{-1}}},\qquad\varepsilon(E)=(-1)^{\unsure{\varepsilon_{-1}|e_{-1}}},
\]
donc
\[\boxed{\varepsilon'(E)=(-1)^{\unsure{\varepsilon_1|e_1}}}.\]

Lorsque l'on n'est pas dans le cas favorable $X$ lisse, $E$ \unsure{c$^{\mathrm{t}}$|ct} \unsure{tordu}, $\deleted{\unsure{E}}\check E$ \uline{globalement} isomorphe à $E(\rho)$, alors $D(E)$ ne s'exprime plus de façon aussi simple en fonction de $E$, et l'équation fonctionnelle (1) ne s'exprime plus \unsure{nec.|avec} en termes de la seule fonction $L_E$ [mais elle relie \uline{deux} fonctions $L_{DE}$, $L_E$]
\mspage{16}
